\documentclass[journal]{IEEEtran}

\usepackage{amsmath}
\usepackage{booktabs}
\usepackage{graphicx}
\usepackage{url}

\title{Leakage-Controlled Evaluation of Feature Stability and Selective Reliability in CARDIODIAG}

\author{[AUTHOR NAMES AND AFFILIATIONS TO BE ADDED]}

\begin{document}
\maketitle

\begin{abstract}
This manuscript reports a leakage-controlled reproduction and extension of the
CARDIODIAG cardiovascular disease prediction framework. The study separates the
historical notebook protocol from a corrected pipeline in which imputation,
encoding, outlier handling, feature selection, model selection, and uncertainty
calibration respect data boundaries. A stability-aware training-only feature
selection component is evaluated against corrected logistic, voting, and
stacking baselines, and a separate selective-reliability layer measures whether
confidence ranks correct decisions above errors. Internal nested-validation
results are generated directly from machine-readable experiment artifacts.
Independent external validation is reported only if a compatible dataset is
frozen before final model selection; otherwise, the absence of independent
validation is treated as a limitation.
\end{abstract}

\begin{IEEEkeywords}
cardiovascular disease, leakage control, nested validation, explainable AI,
calibration, conformal prediction, selective prediction.
\end{IEEEkeywords}

\section{Introduction}
Cardiovascular prediction models must be evaluated for discrimination,
calibration, robustness, and transportability rather than accuracy alone.
CARDIODIAG provides an important starting point because it combines explainable
feature reduction with optimized ensemble classifiers \cite{khan2026cardiodiag}.
This study asks two falsifiable questions under properly separated validation:
whether a resampled feature selector is more stable than the same relevance
score applied once, and whether an abstention layer can identify higher-risk
predictions. The latter is treated as empirical selective reliability rather
than evidence that a model semantically understands uncertainty.

\textbf{Contribution status.} The proposed contribution is a
leakage-controlled stability--reliability decoupling audit for CARDIODIAG, not
a confirmed novel algorithm. It tests whether more stable selected feature sets
actually yield better discrimination, calibration, or selective reliability.
The current evidence is internal and does not establish clinical utility,
deployment readiness, or external validity.

\section{Related Work and Research Gap}
Feature-selection stability is an established property of a selection procedure,
not a substitute for predictive validation. Nogueira \emph{et al.} provide a
chance-corrected stability statistic for comparing selected sets
\cite{nogueira2018stability}, while resampling-based stability selection has
distinct assumptions and error-control goals \cite{meinshausen2010stability}.
Nested validation is required when the same data support model selection and
performance estimation \cite{varma2006bias,cawley2010overfitting}. Calibration,
discrimination, and selective risk are likewise distinct properties of a
prediction model \cite{vancalster2019calibration}. Conformal prediction can
provide marginal coverage under exchangeability, but coverage should not be
assumed under population shift \cite{angelopoulos2023conformal}.

The gap is narrower than a new feature-selection algorithm: this study tests
whether training-only feature-set stability in the CARDIODIAG representation is
associated with held-out predictive and selective reliability. The reporting
and risk-of-bias framing follows TRIPOD+AI and PROBAST+AI
\cite{collins2024tripodai,moons2025probastai}.

\section{Audit of the Original Implementation}
The supplied repository is notebook-centric and contains multiple historical
protocols. The historical reproduction and its deviations are documented in
\texttt{docs/baseline\_reproduction\_report.md}. In particular, available
notebooks contain paths where resampling, preprocessing, and dimensionality
reduction occur before the evaluation split.

\section{Methods}
\subsection{Corrected Pipeline}
All data-dependent operations are fitted on the training data of the relevant
fold. Feature types are declared in the local dataset manifest: nominal and
ordinal coded predictors are one-hot encoded, identifiers are excluded, and
numeric missingness/outlier treatment is training-only. Unknown external data
must provide an explicit feature mapping before the external gate can pass.

\subsection{Stability-Aware Feature Reduction}
For feature $j$ and inner resample $b$, let $r_{j,b}$ denote the combined
training-only relevance score and $s_{j,b}$ its selection indicator. Relevance
is estimated after splitting each resample into a fitting part and a validation
part; permutation importance is evaluated on the latter, which remains inside
the outer training boundary. We define
\begin{equation}
\pi_j = \frac{1}{B}\sum_{b=1}^{B}s_{j,b}, \qquad
\bar r_j = \frac{1}{B}\sum_{b=1}^{B}r_{j,b}.
\end{equation}
Features are retained when their selection frequency and mean relevance pass
the prespecified thresholds. To test whether this improves stability, the
selected sets are compared with a single-pass relevance selector over matched
repeated stratified training perturbations using the chance-corrected stability
statistic of Nogueira \emph{et al.} \cite{nogueira2018stability}. The reported
interval is a paired bootstrap interval over repeat-level stability estimates.
This is an operational stability rule, not a causal interpretation of feature
importance.

\subsection{Uncertainty and Calibration}
The project reports Brier score, reliability-bin error, calibration intercept
and slope, split-conformal prediction-set coverage, abstention rate, confidence-
error AUC, area under the risk-coverage curve (AURC), and selective risk. The
confidence-error AUC tests whether confidence ranks correct predictions above
errors; it is not a probability of correctness. Conformal coverage is
interpreted under its exchangeability assumption \cite{angelopoulos2023conformal}
and is measured again under any available dataset shift.

\section{Experimental Protocol}
The historical protocol, corrected common-split protocol, and nested-validation
protocol are kept separate. Hyperparameter selection is performed only in inner
folds. The final internal holdout is not repeatedly consulted. Reporting follows
the transparency principles of TRIPOD+AI \cite{collins2024tripodai} and the
development/evaluation distinctions in PROBAST+AI \cite{moons2025probastai}.

For feature stability, five stratified folds were repeated ten times. Each
training fold was transformed using training data only, and the stable selector
was compared with a one-shot selector on the same transformed feature universe.
For selective reliability, ten random seeds used a 56.25\% fitting, 18.75\%
calibration, and 25\% evaluation split. No evaluation labels were used for
fitting, calibration, threshold selection, or model selection.

The primary endpoint is paired outer-fold ROC-AUC for the stable candidate
against its matched corrected candidate. PR-AUC, Brier score, calibration,
feature stability, and selective metrics are secondary. The repeated intervals
are descriptive paired bootstrap summaries over the recorded perturbations or
seeds; they are not treated as independent-patient confidence intervals.

\section{Datasets and Provenance}
The supplied Five States and Mendeley datasets are development/replication
benchmarks because they were already used by the base study. Their hashes and
roles are stored in \texttt{data/manifests/local\_datasets.json}. A dataset
used for external validation must be frozen in
\texttt{data/manifests/external\_dataset.json} before any label-based modeling
decision.

\section{Results}
Table~\ref{tab:nested} preserves the earlier corrected nested-validation
artifacts generated before the explicit coded-category schema correction. It is
an internal result only and must not be described as independent external
evidence. The post-correction bounded sensitivity results are reported
separately below so that the earlier artifact is not silently overwritten.

\begin{table*}[t]
\centering
\caption{Nested internal validation summary.}
\label{tab:nested}
\begin{tabular}{llrrrrr}
\toprule
Dataset & Candidate & ROC-AUC & PR-AUC & Brier & ECE & ROC-AUC CI \\
\midrule
heart5.csv & corrected\_logistic & 0.928 $\pm$ 0.006 & 0.935 & 0.107 & 0.104 & [0.907, 0.942] \\
heart5.csv & corrected\_ovec & 0.930 $\pm$ 0.007 & 0.934 & 0.100 & 0.109 & [0.909, 0.944] \\
heart5.csv & corrected\_osec & 0.926 $\pm$ 0.010 & 0.929 & 0.109 & 0.156 & [0.899, 0.936] \\
heart5.csv & stable\_logistic & 0.928 $\pm$ 0.006 & 0.935 & 0.107 & 0.101 & [0.907, 0.942] \\
heart5.csv & stable\_osec & 0.925 $\pm$ 0.011 & 0.929 & 0.109 & 0.156 & [0.896, 0.934] \\
MendeleyCardiovascular\_Disease\_Dataset.csv & corrected\_logistic & 0.994 $\pm$ 0.003 & 0.995 & 0.028 & 0.172 & [0.988, 0.996] \\
MendeleyCardiovascular\_Disease\_Dataset.csv & corrected\_ovec & 0.996 $\pm$ 0.001 & 0.997 & 0.026 & 0.136 & [0.993, 0.998] \\
MendeleyCardiovascular\_Disease\_Dataset.csv & corrected\_osec & 0.995 $\pm$ 0.002 & 0.996 & 0.024 & 0.211 & [0.992, 0.998] \\
MendeleyCardiovascular\_Disease\_Dataset.csv & stable\_logistic & 0.994 $\pm$ 0.003 & 0.995 & 0.028 & 0.188 & [0.989, 0.996] \\
MendeleyCardiovascular\_Disease\_Dataset.csv & stable\_osec & 0.996 $\pm$ 0.002 & 0.997 & 0.024 & 0.216 & [0.990, 0.997] \\
\bottomrule
\end{tabular}

\end{table*}

\subsection{Schema-Corrected Sensitivity Rerun}
The schema-corrected bounded rerun used two outer and two inner folds for each
candidate. The results were generated directly from
\texttt{results/generated/schema\_corrected\_*nested\_2x2.json}; the source
table is reproduced in Table~\ref{tab:schema_corrected}. Corrected OVEC had the
highest bounded ROC-AUC mean on both datasets (0.928 on Five States and 0.995 on
Mendeley). Stable logistic was similar on Five States, while stable OSEC had
higher ECE than corrected OVEC on both datasets. Because this rerun is bounded,
it is sensitivity evidence rather than a replacement for a full confirmatory
analysis.

\begin{table*}[t]
\centering
\caption{Schema-corrected bounded nested sensitivity results.}
\label{tab:schema_corrected}
\begin{tabular}{llrrrrr}
\toprule
Dataset & Candidate & ROC-AUC & PR-AUC & Brier & ECE & ROC-AUC CI \\
\midrule
heart5.csv & corrected\_logistic & 0.927 $\pm$ 0.000 & 0.932 & 0.109 & 0.087 & [0.904, 0.939] \\
heart5.csv & corrected\_ovec & 0.928 $\pm$ 0.000 & 0.929 & 0.101 & 0.090 & [0.906, 0.942] \\
heart5.csv & corrected\_osec & 0.925 $\pm$ 0.000 & 0.929 & 0.113 & 0.103 & [0.892, 0.931] \\
heart5.csv & stable\_logistic & 0.927 $\pm$ 0.000 & 0.934 & 0.108 & 0.088 & [0.905, 0.940] \\
heart5.csv & stable\_osec & 0.925 $\pm$ 0.003 & 0.930 & 0.114 & 0.122 & [0.903, 0.939] \\
MendeleyCardiovascular\_Disease\_Dataset.csv & corrected\_logistic & 0.989 $\pm$ 0.004 & 0.991 & 0.032 & 0.133 & [0.980, 0.993] \\
MendeleyCardiovascular\_Disease\_Dataset.csv & corrected\_ovec & 0.995 $\pm$ 0.002 & 0.996 & 0.030 & 0.107 & [0.991, 0.997] \\
MendeleyCardiovascular\_Disease\_Dataset.csv & corrected\_osec & 0.992 $\pm$ 0.003 & 0.994 & 0.032 & 0.212 & [0.978, 0.992] \\
MendeleyCardiovascular\_Disease\_Dataset.csv & stable\_logistic & 0.990 $\pm$ 0.005 & 0.992 & 0.032 & 0.146 & [0.984, 0.994] \\
MendeleyCardiovascular\_Disease\_Dataset.csv & stable\_osec & 0.993 $\pm$ 0.004 & 0.994 & 0.028 & 0.238 & [0.988, 0.996] \\
\bottomrule
\end{tabular}

\end{table*}

\subsection{Reliability}
Reliability artifacts are stored under \texttt{results/generated/}. The
repeated internal analysis is summarized in Table~\ref{tab:uncertainty}. The
confidence-error AUC exceeded 0.5 for both datasets, and conformal coverage was
close to the nominal 90\% target. The stability-aware selector did not improve
all reliability quantities consistently.

\begin{table*}[t]
\centering
\caption{Repeated internal selective-reliability evidence. Values are means
over ten seeds; confidence intervals are stored in the machine-readable result
files.}
\label{tab:uncertainty}
\begin{tabular}{llrrrr}
\toprule
Dataset & Candidate & Conf.-error AUC & AURC & Conformal coverage & ECE \\
\midrule
Five States & Corrected OSEC & 0.762 & 0.0505 & 0.914 & 0.122 \\
Five States & Stable OSEC & 0.760 & 0.0516 & 0.915 & 0.121 \\
Mendeley & Corrected OSEC & 0.905 & 0.0041 & 0.914 & 0.221 \\
Mendeley & Stable OSEC & 0.905 & 0.0042 & 0.914 & 0.217 \\
\bottomrule
\end{tabular}

\end{table*}

\subsection{Feature-Selection Stability}
The stability comparison is shown in Table~\ref{tab:stability}. The stable
selector was more stable than the one-shot comparator on the Mendeley benchmark
with a paired difference of 0.078 and 95\% bootstrap interval [0.035, 0.119].
On the Five States benchmark, the difference was 0.098 but the interval
[-0.111, 0.254] included zero. Thus, the evidence supports a dataset-conditional
stability effect, not a universal claim.

\begin{table}[t]
\centering
\caption{Repeated feature-selection stability evidence. $\Phi$ is the
chance-corrected stability statistic.}
\label{tab:stability}
\begin{tabular}{lrrrr}
\toprule
Dataset & Stable $\Phi$ & One-shot $\Phi$ & Difference & 95\% bootstrap CI \\
\midrule
Five States & 0.112 & 0.014 & 0.097 & [-0.027, 0.198] \\
Mendeley & 0.284 & 0.215 & 0.069 & [0.000, 0.167] \\
\bottomrule
\end{tabular}

\end{table}

\subsection{Stability--Reliability Linkage}
Table~\ref{tab:linkage} tests whether a more stable selected set translated into
better held-out reliability on the same repeated splits. Differences are
stable OSEC minus corrected OSEC; positive AURC or ECE differences indicate
worse selective risk or calibration, respectively. The paired intervals include
zero for the main discrimination and reliability comparisons. On Five States,
the stable candidate had a small adverse AURC difference; on Mendeley, the
differences were close to zero. Thus, selector stability did not act as a
surrogate for predictive or selective reliability.

\begin{table*}[t]
\centering
\caption{Paired stability--reliability linkage. $\Phi$ is cross-seed selected-
set stability. $\Delta$ values are stable OSEC minus corrected OSEC means; the
full bootstrap intervals are stored in the machine-readable linkage artifacts.}
\label{tab:linkage}
\begin{tabular}{lrrrrr}
\toprule
Dataset & Cross-seed $\Phi$ & $\Delta$ ROC-AUC & $\Delta$ Conf.-AUC & $\Delta$ AURC & $\Delta$ ECE \\
\midrule
Five States & 0.201 & -0.0009 & -0.0015 & +0.0011 & -0.0001 \\
Mendeley & 0.178 & -0.0002 & -0.0001 & +0.0002 & -0.0045 \\
\bottomrule
\end{tabular}

\end{table*}

\section{Discussion}
The nested results do not support a uniform performance advantage for the
stability-aware selector. On the Five States benchmark, corrected OVEC has the
highest mean ROC-AUC among the listed candidates, while stable OSEC is lower.
On the Mendeley benchmark, corrected OVEC and stable OSEC are similar. The
schema-corrected stability comparison shows positive paired differences on both
supplied benchmarks, but the stability intervals are based on only three
repeat-level comparisons. The four-seed reliability results do not show
consistent improvement across discrimination, selective risk, coverage, and
calibration. These findings separate feature-set stability from predictive
reliability: more repeatable feature selection does not by itself establish
better calibrated or safer predictions. A selector is not justified merely
because it reduces dimensionality or includes an explainability component.

\section{External Validation and Threats to Validity}
No independent compatible dataset was available for a defensible frozen
external evaluation within this project. The supplied Mendeley dataset is not
new external evidence because it was used in the base paper, and the public UCI
component sources overlap the Five States composite. Consequently, this study
reports internal nested validation and reliability analysis only, and does not
claim independent external validation. This caution is consistent with the
documented heterogeneity and performance degradation commonly observed when
cardiovascular prediction models are evaluated outside their development data
\cite{wessler2021external}.

\section{Statistical Analysis and Robustness}
The study prespecifies ROC-AUC as the primary comparison and reports PR-AUC,
classification metrics, Brier score, calibration intercept and slope, uniform
ten-bin ECE, selected-set stability, confidence-error AUC, AURC, and conformal
coverage as secondary quantities. Paired comparisons use identical partitions
for the candidates. No multiplicity-adjusted confirmatory claim is made across
the secondary metrics. Synthetic missingness or measurement perturbations,
when run, will be reported as stress tests and not as independent validation.

The bounded post-correction stability runs found positive comparative stability
for the stable selector on both local benchmarks: the paired difference was
0.189 [0.027, 0.425] on Five States and 0.269 [0.039, 0.574] on Mendeley.
However, the four-seed linkage runs found no consistent reliability advantage:
the paired ROC-AUC difference was 0.0009 [0.0002, 0.0019] on Five States and
-0.0004 [-0.0013, 0.0004] on Mendeley, while calibration and selective metrics
did not uniformly favor the stable candidate. These are bounded internal
results and remain subject to dependence and small-repeat uncertainty.

An evaluation-only stress test at seed 42 applied 10\% independent missingness
and numeric noise with standard deviation equal to 5\% of each training feature
standard deviation. On Mendeley, ROC-AUC fell from 0.998 to 0.810 under the
numeric-noise perturbation for corrected OSEC and from 0.998 to 0.815 for stable
OSEC. On Five States, 10\% missingness reduced ROC-AUC from approximately 0.935
to 0.922 for both candidates. These are synthetic single-split stress tests,
not independent external validation.

\section{Limitations and Ethical Considerations}
The study is retrospective, benchmark-based, and not a clinical diagnostic
device. It does not establish prospective utility, patient benefit, fairness
across unsupported demographic groups, or clinical safety. Dataset provenance,
license conditions, outcome definitions, and possible source overlap constrain
transportability.

\section{Reproducibility}
The code, tests, dependency lock, data manifests, protocol reports, and result
generators are maintained in the repository. Raw restricted data must not be
committed.

\section{Conclusion}
This work delivers a reproducible leakage-controlled audit and a falsifiable
stability--reliability decoupling audit for CARDIODIAG. Internal evidence shows
that feature-set stability and selective reliability can move independently:
the stable selector does not produce a consistent predictive or reliability
benefit despite positive comparative stability evidence on one benchmark.
Independent external and prospective evaluation remain necessary before any
generalization or clinical-readiness claim.

\bibliographystyle{IEEEtran}
\bibliography{references}
\end{document}
